\section{Introduction}
\label{sec:intro}

Zone-graph reachability appears tailor-made for a GPU. A BFS level contains millions of
$(\text{state},\text{edge})$ pairs; each is expanded independently; the arithmetic is small
dense integer matrix work. Yet across a sustained implementation effort we were unable to
make a GPU implementation beat a 12-core CPU implementation of the same algorithm on the same
chip, and the reasons turned out to be interesting enough to report.

This paper is a negative result with a mechanism. Its value is threefold. First, it refutes,
\emph{on measurement rather than intuition}, the hypothesis most practitioners (including the
authors) reach for first --- branch divergence --- and thereby saves the reader from
implementing the standard fix. Second, it identifies the actual cause and shows it is not a
coding failure but a consequence of the algorithm's shape. Third, it delimits the claim
precisely: the dominant cause is a hardware capacity that a datacenter GPU relaxes, so the
result should not be read as ``GPUs cannot do this''.

\paragraph{Contributions.}
\begin{itemize}
\item Three GPU DBM representations, all implemented and validated against exact
\uppaal{} state counts (Sect.~\ref{sec:reps}).
\item An adaptive per-level CPU/GPU merge router, whose selections are themselves a
measurement instrument (Sect.~\ref{sec:adaptive}).
\item A refutation of the divergence hypothesis and identification of the working-set
dilemma (Sect.~\ref{sec:why}), plus a counting argument showing no representation escapes it
(Sect.~\ref{sec:cooperation}).
\item A scoped outlook for datacenter GPUs (Sect.~\ref{sec:h100}).
\end{itemize}
